\section{Conclusion}
\label{sec:conclusion}

Autonomous storage maintenance is not solved by attaching a failure predictor to a migration command. Predictions are uncertain, maintenance actions interact, background repair consumes foreground resources, and an automated transition can be harmful even when the underlying repair primitive is correct. This paper presents \autopilot{}, a file-system-local controller that places persistent evidence and an explicit safety envelope between observations and mutations, then executes scrub, rebalance, and proactive drain as bounded, auditable work.

The \argosfs{} implementation demonstrates confirmation and cooldown hysteresis, drain survivability checks, conflict-aware fail-closed mutation, resumable scrub/rebalance, foreground-pressure budgeting, outcome history, post-action verification, and dry-run explanations. It also exposes clear current limits: the controller command targets host volumes, foreground rate control is coarse, and several reserved policy fields are not yet active.

\resulttext{insert the final one-sentence quantitative conclusion on safety, risk exposure, and foreground interference.} Regardless of the final effect sizes, the evaluation is structured so the central claim is falsifiable against a simple fixed-threshold controller receiving the same health evidence. The resulting question is not whether storage can be made ``self-driving'' in the abstract, but whether a closed-loop controller can earn the authority to mutate a resilient file system by demonstrating conservative decisions under imperfect information.
